% ---------------------------------------------------------------------------
% formal_framework.tex
% Sidecar (Fable, review round 2). NOT inputted by main.tex.
% Harvested 7 Sep 2026 into ch 3 (Markov game, Bellman), ch 4 (operators,
% logit-update, shaper objective), ch 6 (estimators, H1--H3).
% Fully wired 7 Sep 2026 (round 3): every labelled equation below now
% lives in chapters 3, 4 or 6 with the same label, plus tab:hyperparameters.
% Corrections made while wiring (do not copy the text below verbatim):
%   * sigma_B is 13-15 per update (live_adv logs), not 53; and under Adam
%     the scale is not a step size (ch 4, "What the normalisation operator
%     changes").
%   * [CONFIRM] items resolved from TRL 0.11.4: naive clip 0.2; KL
%     controller adaptive, init 0.2, target 6, horizon 1e4 (logged
%     0.200 -> 0.181 over 75 updates).
%   * Joint optimum under xi is 207.05 unconstrained; 206.6 is the
%     symmetric optimum. cpr_xi.py now prints both and the BR Q-vector.
%   * H2 uses leave-2 at R<12, not pi(1|R<12), as pre-registered.
%   * pi(1) at the opening is about 0.10 (first-episode records), not 0.05.
% Not harvested as-is: Perolat metrics as reported quantities,
% pi(1|R<12) as H2, scripts/noise_dp.py (live source is cpr_xi.py),
% unconfirmed KL coefficients, Barrett 2012.
% ---------------------------------------------------------------------------
% Three drop-in sections that define every object the thesis measures.
%   Section 1 -> Chapter 3 (Environment), before the constant-pair table.
%   Section 2 -> Chapter 4 (Method), replacing the "LoRA and PPO",
%                "Advantage normalisation" and "Two-learner / shaper path"
%                prose with these definitions plus your hyperparameter table.
%   Section 3 -> Chapter 6 (Experimental design), before the executed runs.
%
% Requires amsmath, amssymb, booktabs (already loaded by main.tex).
% Labels defined here: eq:harvest, eq:regen, eq:growth, eq:return,
%   eq:collapse, eq:metrics, eq:bellman-br, eq:q-open, eq:bellman-joint,
%   eq:policy, eq:klreward, eq:gae, eq:norm, eq:ppo, eq:logit-update,
%   eq:shaper-objective, eq:naive-update, eq:estimators, eq:wilson,
%   eq:paired, eq:h1, eq:h2, eq:h3.
% Numbers quoted in the worked examples are your logged A0 values
% (39.8 / 6.2 raw, 1.38 / 0.75 whitened, 16.7 / -5.7 centred) and the
% untrained prior (about 0.90 on token 2). Replace pi(1) = 0.05 with the
% measured pre-flight value.
% Two items to confirm against the config before submission are marked
% [CONFIRM].
% ---------------------------------------------------------------------------

% ===========================================================================
\section{Problem formulation}
\label{sec:formulation}
% ===========================================================================

\paragraph{The game.}
The environment is a two-player, finite-horizon Markov game
\(\mathcal{G} = (\mathcal{S}, \mathcal{A}, P, c, T, R_0)\).
The state is the pre-harvest stock
\(R_t \in \mathcal{S} = \{0, 1, \dots, K\}\) at the start of round
\(t \in \{1, \dots, T\}\), with \(R_1 = R_0\).
Each player \(i \in \{1, 2\}\) simultaneously submits a request
\(a_{i,t} \in \mathcal{A} = \{0, 1, 2, 3\}\).
Throughout, \(K = 40\), \(R_0 = 8\), \(T = 36\), and the growth rate is
\(\rho = 0.9\).

\paragraph{Harvest.}
Requests are filled in full when the pool can cover them and are capped
otherwise.
Writing \(A_t = a_{1,t} + a_{2,t}\), the receipts \(c_{i,t}\) and the
post-harvest stock \(\tilde{R}_t\) are
\begin{equation}
  (c_{1,t}, c_{2,t}, \tilde{R}_t) =
  \begin{cases}
    \bigl(a_{1,t},\; a_{2,t},\; R_t - A_t\bigr) & \text{if } A_t \le R_t, \\[2pt]
    \bigl(\min(a_{1,t}, \lfloor R_t/2 \rfloor),\;
          \min(a_{2,t}, \lfloor R_t/2 \rfloor),\; 0\bigr) & \text{if } A_t > R_t .
  \end{cases}
  \label{eq:harvest}
\end{equation}
The second branch is scarcity: whatever is not received is wasted and
the pool is emptied.

\paragraph{Regeneration.}
A strictly positive post-harvest stock regrows by a rounded logistic
increment, scaled by a multiplier \(\xi_t\); an empty pool is absorbing:
\begin{align}
  R_{t+1} &=
  \begin{cases}
    \min\bigl(K,\; \tilde{R}_t + g(\tilde{R}_t; \xi_t)\bigr) & \tilde{R}_t > 0, \\
    0 & \tilde{R}_t = 0,
  \end{cases}
  \label{eq:regen} \\[4pt]
  g(R; \xi) &= \Bigl\lfloor \xi\, \rho\, \frac{R\,(K - R)}{K} \Bigr\rceil ,
  \label{eq:growth}
\end{align}
where \(\lfloor \cdot \rceil\) denotes rounding half to even, evaluated on
the exact rational argument.
In the deterministic game (Stage~A) \(\xi_t \equiv 1\).
In the stochastic game (Stage~B) the multipliers are independent and
identically distributed, uniform on \(\{0.7, 1.0, 1.3\}\), independent of
all actions, and common to both players within a game.
Setting \(\xi_t \equiv 1\) in \eqref{eq:growth} recovers the deterministic
update exactly, so Stage~A is the special case of Stage~B.

\paragraph{Return, collapse, survival.}
Player \(i\)'s return is the undiscounted sum of receipts,
\begin{equation}
  G_i = \sum_{t=1}^{T} c_{i,t} .
  \label{eq:return}
\end{equation}
The collapse round is the first round after which the pool is empty,
\begin{equation}
  \tau = \min\{\, t : \tilde{R}_t = 0 \,\},
  \qquad
  S = \mathbb{1}[\tau > T],
  \label{eq:collapse}
\end{equation}
with \(\tau = \infty\) if the pool never empties, so \(S = 1\) means the
pool survived the horizon.
Because \(g \ge 0\) and \(\xi_t \ge 0\), the stock can fall only through
harvest; collapse is therefore always the consequence of the players'
requests.

\paragraph{Social metrics.}
For any pair of policies the outcome is summarised by three quantities
in the sense of P\'erolat et al.\ \cite{perolat2017}:
\begin{equation}
  \begin{aligned}
  \text{sustainability} &= \Pr(S = 1), &
  \text{efficiency} &= \frac{\mathbb{E}[G_1 + G_2]}{W^{*}}, \\[2pt]
  \text{equality} &= 1 - \frac{\mathbb{E}\,|G_1 - G_2|}{\mathbb{E}[G_1 + G_2]},
  \end{aligned}
  \label{eq:metrics}
\end{equation}
where \(W^{*}\) is the joint optimum defined in \eqref{eq:bellman-joint}
(\(210\) in Stage~A, \(206.6\) in Stage~B).

\paragraph{Policies and the constant-strategy reduction.}
A policy \(\pi_i\) maps the player's observable history to a distribution
over \(\mathcal{A}\).
The constant policies \(\pi_i \equiv k\), \(k \in \mathcal{A}\), induce a
\(4 \times 4\) matrix of returns and collapse rounds,
\(M(k_1, k_2) = (G_1, G_2, \tau)\), which is the empirical payoff matrix of
Leibo et al.\ \cite{leibo2017} restricted to constant strategies.
Table~\ref{tab:chicken} reports its relevant cells.
Under this reduction the game is a chicken game (the asymmetric profiles
\((2,1)\) and \((1,2)\) are the strict pure equilibria and \((1,1)\) is not
an equilibrium), with the degenerate feature that the dove's return is the
same, \(36\), whether or not it is exploited.

\paragraph{Best response to a fixed opponent.}
Against an opponent playing a fixed policy \(\sigma\), the optimal value of
player~1 satisfies the Bellman equation
\begin{equation}
  V^{\sigma}(R, t) = \max_{a \in \mathcal{A}}
  \Bigl\{ c_1\bigl(R, a, \sigma(R, t)\bigr)
  + \mathbb{E}_{\xi}\bigl[ V^{\sigma}(R', t+1) \bigr] \Bigr\},
  \qquad V^{\sigma}(\cdot, T+1) = 0,
  \label{eq:bellman-br}
\end{equation}
with \(R'\) given by \eqref{eq:harvest}--\eqref{eq:growth}.
The value of each opening action is
\begin{equation}
  Q^{\sigma}_{\mathrm{open}}(a) = c_1\bigl(R_0, a, \sigma(R_0, 1)\bigr)
  + \mathbb{E}_{\xi}\bigl[ V^{\sigma}(R', 2) \bigr] .
  \label{eq:q-open}
\end{equation}
Against the frozen hawk \(\sigma \equiv 2\), \eqref{eq:q-open} gives
\(Q_{\mathrm{open}} = (105, 104, 102, 6)\) in Stage~A and
\((102.5, 101.7, 99.1, 91.1)\) in Stage~B for openings \(0, 1, 2, 3\).

\paragraph{Joint optimum.}
The largest achievable joint return is
\begin{equation}
  W^{*}(R, t) = \max_{(a_1, a_2) \in \mathcal{A}^2}
  \Bigl\{ c_1 + c_2 + \mathbb{E}_{\xi}\bigl[ W^{*}(R', t+1) \bigr] \Bigr\},
  \qquad W^{*} := W^{*}(R_0, 1),
  \label{eq:bellman-joint}
\end{equation}
attained in Stage~A by one round of \((0,0)\) followed by \((3,3)\) at the
fixed point \(R = 14\), for \(105\) per player.
Since the state space has \(41\) stocks, \(36\) rounds and at most three
draws, \eqref{eq:bellman-br} and \eqref{eq:bellman-joint} are solved
exactly by backward induction in rational arithmetic
(\texttt{scripts/noise\_dp.py}); every dynamic-programming value quoted
in this thesis is an output of that computation.

% ===========================================================================
\section{Learning algorithm}
\label{sec:algorithm}
% ===========================================================================

\paragraph{Policy.}
Each learning agent is Gemma-2-2b-it with a rank-2 LoRA adapter
\(\theta\) on the query and value projections and a scalar value head.
At round \(t\) the agent receives a prompt \(x_t\) rendered from its
observable history (the current stock, the previous round's requests and
receipts, and, for a shaper with the trial prompt enabled, the joint
request counts and episode summaries of the current trial).
The action is one token drawn from the softmax over the four legal
action tokens, all other vocabulary entries being masked:
\begin{equation}
  \pi_{\theta}(a \mid x_t) = \frac{\exp z_{\theta}(x_t)_a}
  {\sum_{a' \in \mathcal{A}} \exp z_{\theta}(x_t)_{a'}},
  \qquad a \in \mathcal{A},
  \label{eq:policy}
\end{equation}
where \(z_{\theta}(x_t)\) are the model's logits at the answer position.
The untrained model places about \(0.90\) of its mass on token~2 at the
opening prompt.

\paragraph{Per-step reward.}
The reward handed to the optimiser is the receipt, penalised by the
per-token divergence from the frozen reference model (the same weights
with the adapter disabled):
\begin{equation}
  \tilde{r}_t = c_t - \beta_t \bigl[\log \pi_{\theta}(a_t \mid x_t)
  - \log \pi_{\mathrm{ref}}(a_t \mid x_t)\bigr],
  \label{eq:klreward}
\end{equation}
with \(\beta_t\) set by the adaptive controller of the TRL implementation
[CONFIRM: initial coefficient and target from the config].
Steps after collapse (\(R_t = 0\)) are masked out of every batch
statistic and out of the loss; the collapse step itself is not masked.

\paragraph{Advantages.}
Advantages are generalised advantage estimates over each episode
(naive agents) or over the concatenated episodes of a trial (shapers):
\begin{equation}
  \delta_t = \tilde{r}_t + \gamma V_{\theta}(x_{t+1}) - V_{\theta}(x_t),
  \qquad
  \hat{A}_t = \sum_{l \ge 0} (\gamma \lambda)^{l}\, \delta_{t+l},
  \qquad \gamma = 1,\ \lambda = 0.97 .
  \label{eq:gae}
\end{equation}
Before entering the loss, the advantages of a batch \(\mathcal{B}\) are
transformed by one of two operators, where \(\mu_{\mathcal{B}}\) and
\(\sigma_{\mathcal{B}}\) are the mean and standard deviation over the
unmasked positions of the batch:
\begin{equation}
  \mathcal{N}_{\mathrm{whiten}}(\hat{A}_t) = \frac{\hat{A}_t - \mu_{\mathcal{B}}}
  {\sqrt{\sigma_{\mathcal{B}}^2 + \varepsilon}},
  \qquad
  \mathcal{N}_{\mathrm{centre}}(\hat{A}_t) = \hat{A}_t - \mu_{\mathcal{B}} .
  \label{eq:norm}
\end{equation}
Whitening is the TRL default (\texttt{masked\_whiten}); centring is the
variant used in the centred runs.
The two operators differ only by the factor \(1/\sigma_{\mathcal{B}}\), so
at a fixed learning rate every centred update is \(\sigma_{\mathcal{B}}\)
times larger than the corresponding whitened one.

\paragraph{Objective.}
With \(\tilde{A}_t = \mathcal{N}(\hat{A}_t)\), the probability ratio
\(\varrho_t = \pi_{\theta}(a_t \mid x_t) / \pi_{\theta_{\mathrm{old}}}(a_t \mid x_t)\),
and the empirical return-to-go \(\hat{G}_t\), each agent maximises the
clipped surrogate
\begin{equation}
  \begin{aligned}
  \mathcal{L}(\theta) = \mathbb{E}_{\mathcal{B}}\Bigl[
    &\min\bigl(\varrho_t \tilde{A}_t,\;
    \operatorname{clip}(\varrho_t, 1 - \epsilon, 1 + \epsilon)\, \tilde{A}_t\bigr) \\
    &- c_v \bigl(V_{\theta}(x_t) - \hat{G}_t\bigr)^2
    + c_e\, \mathcal{H}\bigl(\pi_{\theta}(\cdot \mid x_t)\bigr)
  \Bigr],
  \end{aligned}
  \label{eq:ppo}
\end{equation}
with \(\epsilon = 0.2\) for naive learners and \(0.1\) for shapers
[CONFIRM the naive value], \(c_v = 0.01\), and \(c_e\) decaying from
\(0.05\) to \(0.01\) over the first 150 updates; one pass over the batch
in minibatches of five; Adam with learning rate \(1.41 \times 10^{-6}\)
(naive) or \(3 \times 10^{-7}\) (shaper).
The remaining hyperparameters are collected in
Table~\ref{tab:hyperparameters}.
% [Add a single table with every hyperparameter: LR, clip, batch,
%  minibatch, PPO epochs, gamma, lambda, vf_coef, entropy schedule, KL
%  coefficient and controller, LoRA rank/alpha/dropout/targets, seeds,
%  epochs, games per epoch, episodes per trial, hardware, wall-clock.]

\paragraph{Expected update on the opening logits.}
The effect of the normalisation operator on the opening decision can be
read from the surrogate at \(\varrho_t = 1\).
For a softmax policy, \(\partial \log \pi(a \mid x) / \partial z_k
= \mathbb{1}[a = k] - \pi(k \mid x)\), so the expected gradient of the
surrogate with respect to the logit of opening action \(k\) is
\begin{equation}
  \frac{\partial \mathcal{L}}{\partial z_k}\Big|_{x_1}
  = \pi(k \mid x_1)\bigl(\tilde{A}(k) - \bar{A}\bigr),
  \qquad
  \bar{A} = \sum_{a \in \mathcal{A}} \pi(a \mid x_1)\, \tilde{A}(a),
  \label{eq:logit-update}
\end{equation}
where \(\tilde{A}(a)\) is the normalised opening advantage of action
\(a\).
Two consequences follow.
The sign of the update on logit \(k\) is the sign of
\(\tilde{A}(k) - \bar{A}\), so an action whose advantage is positive but
below the policy-weighted mean is pushed down, not up; and the magnitude
is proportional to \(\pi(k \mid x_1)\), so a rare action moves slowly
however large its advantage.
With the untrained prior \(\pi(2) \approx 0.90\), \(\pi(1) \approx 0.05\)
and the logged opening advantages of Test~A, whitening
(\(\tilde{A}(1) = 1.38\), \(\tilde{A}(2) = 0.75\), \(\bar{A} \approx 0.78\))
gives updates of about \(+0.03\) on logit~1 and \(-0.03\) on logit~2,
whereas centring (\(\tilde{A}(1) = 16.7\), \(\tilde{A}(2) = -5.7\),
\(\bar{A} \approx -4.3\)) gives about \(+1.0\) and \(-1.3\): the same
signs, but a step some forty times larger at the same learning rate.
The ratio is \(\sigma_{\mathcal{B}}\), the batch standard deviation of the
raw advantages, which the logged values imply to be about \(53\); it is
inflated by the large raw estimates on the later steps of surviving
episodes that a lightly trained critic (\(c_v = 0.01\)) leaves in
\eqref{eq:gae}.

\paragraph{Naive and shaping objectives.}
Training proceeds in trials of \(E = 5\) episodes.
A naive learner updates its parameters after every episode \(e\) from
that episode's batch,
\begin{equation}
  \theta_{e+1} = \theta_e + \alpha_{\mathrm{n}}\, \nabla_{\theta}
  \mathcal{L}\bigl(\theta_e;\ \tau_e(\pi_{\theta_e}, \pi_{\phi})\bigr),
  \label{eq:naive-update}
\end{equation}
where \(\tau_e\) is the trajectory of episode \(e\) and \(\pi_{\phi}\) is
the partner.
A shaper holds \(\phi\) fixed for the whole trial and updates once, from
the concatenated trajectories, towards
\begin{equation}
  J(\phi) = \mathbb{E}\Bigl[\, \sum_{e=1}^{E} G_{\mathrm{shaper}}^{(e)}
  \bigl(\pi_{\phi}, \pi_{\theta_e}\bigr) \Bigr],
  \qquad \theta_{e+1} = \theta_e + \alpha_{\mathrm{n}} \nabla_{\theta}
  \mathcal{L}(\theta_e; \tau_e(\pi_{\theta_e}, \pi_{\phi})),
  \label{eq:shaper-objective}
\end{equation}
which is the model-free opponent-shaping objective of Lu et al.\
\cite{lu2022} as instantiated for language-model agents by Garcia Segura
et al.\ \cite{garciasegura2025}.
The naive objective in \eqref{eq:naive-update} treats the partner as
fixed within the episode; the shaper objective in
\eqref{eq:shaper-objective} credits \(\phi\) for the partner's parameter
change \(\theta_{e} \to \theta_{e+1}\) that \(\phi\) induced.
That dependence is the only thing that distinguishes shaping from
ordinary learning against a slowly changing opponent, and it can only be
exploited if the shaper's own policy changes on a timescale at which the
naive learner responds.
The shaper's learning rate \(\alpha_{\mathrm{s}} = 3 \times 10^{-7}\),
tighter clip and once-per-trial schedule therefore make it, as a side
effect, a slower and more stationary opponent than a naive learner; the
slow-learning-rate naive--naive control of Section~\ref{sec:executed-two}
isolates that side effect from the shaping term.

% ===========================================================================
\section{Estimators and hypotheses}
\label{sec:estimators}
% ===========================================================================

\paragraph{Estimators.}
An epoch consists of \(n = 15\) games.
For agent \(i\), the opening statistic, survival, and the stock-conditioned
action frequency in an epoch are estimated by
\begin{equation}
  \begin{aligned}
  \hat{L}_i &= \frac{1}{n} \sum_{j=1}^{n} \mathbb{1}\bigl[a^{(j)}_{i,1} \ne 2\bigr],
  &\qquad
  \hat{S} &= \frac{1}{n} \sum_{j=1}^{n} S^{(j)}, \\[2pt]
  \hat{\pi}_i(a \mid R) &= \frac{\sum_{j,t} \mathbb{1}\bigl[R^{(j)}_t = R,\ a^{(j)}_{i,t} = a\bigr]}
  {\sum_{j,t} \mathbb{1}\bigl[R^{(j)}_t = R\bigr]},
  \end{aligned}
  \label{eq:estimators}
\end{equation}
together with the mean return \(\hat{G}_i = \frac{1}{n}\sum_j G_i^{(j)}\).
Last-epoch quantities are reported with the pooled estimate over the last
three epochs (\(n = 45\)) alongside.
For a count \(k\) out of \(n\), the \(95\%\) Wilson interval is
\begin{equation}
  \frac{\hat{p} + z^2/2n \;\pm\; z\sqrt{\hat{p}(1-\hat{p})/n + z^2/4n^2}}{1 + z^2/n},
  \qquad \hat{p} = k/n,\ z = 1.96 ;
  \label{eq:wilson}
\end{equation}
for example \(13/15\) gives \([0.62, 0.96]\) and \(9/15\) gives
\([0.36, 0.80]\).
Seeds are the unit of replication: every contrast between configurations
is reported as the per-seed difference and its mean over the three seeds.
In Stage~B the noise multipliers at a given (epoch, game, round) are
identical across configurations for the same seed, so a between-arm
difference on seed \(s\),
\begin{equation}
  \Delta_s = \hat{Y}^{(\mathrm{arm}\,1)}_s - \hat{Y}^{(\mathrm{arm}\,2)}_s ,
  \label{eq:paired}
\end{equation}
is a paired comparison in which the environmental draws cancel; with three
seeds, agreement of the sign of \(\Delta_s\) across all seeds is reported
rather than a test statistic.

\paragraph{Hypotheses.}
The hypotheses are stated as inequalities between estimated quantities
and values fixed in advance by \eqref{eq:bellman-br}--\eqref{eq:q-open}.
\begin{align}
  \text{H1:}\quad
  & \hat{L}^{\mathrm{B}}_{\mathrm{last\,3}} \in
    \mathrm{CI}_{95}\bigl(\hat{L}^{\mathrm{A}}_{\mathrm{last\,3}}\bigr)
    \ \text{ on each seed (Test~A)},
  \label{eq:h1} \\[3pt]
  \text{H2:}\quad
  & \widehat{\Pr}\bigl(S = 1 \,\big|\, a_{1,1} \ne 2\bigr)
    \ \text{ against } 0.40 \text{ and } 1.00, \quad
    \hat{\pi}_1(1 \mid R < 12) \ \text{ against } 1,
  \label{eq:h2} \\[3pt]
  \text{H3:}\quad
  & \hat{G}_{\mathrm{shaper}} > \hat{G}^{\mathrm{ctrl}}_{\mathrm{hawk}}
    \ \text{ and }\
    \hat{L}^{\mathrm{vs\,shaper}}_{\mathrm{naive}} > \hat{L}^{\mathrm{ctrl}}_{\mathrm{naive}}
    \ \text{ (every seed)},
  \label{eq:h3}
\end{align}
where H1 (opening under variance) asks whether the Stage~B opening share
lies inside the Stage~A interval, H2 (continuation) compares survival
after a restrained opening with the two dynamic-programming values for a
constant continuation and for the feedback rule, and H3 (shaping) asks
whether the shaper out-earns the hawk role of the control while its
partner restrains more.
The control in \eqref{eq:h3} is the slow-learning-rate naive--naive
configuration, so that the comparison holds the partner's learning
timescale fixed and varies only the trial-level objective
\eqref{eq:shaper-objective}.
Failure of either inequality in \eqref{eq:h3} on any seed is reported as
no detectable shaping effect at three seeds.
